% No numbered results (statements-of-record rows resolved here: none).
\section{Introduction}
\label{sec:intro}

When must an obstruction attached to a symmetry vanish?  Here is the
form of the question studied in this paper.  A space of objects \(X\)
carries an involution \(J\), so that applying \(J\) twice returns the
original object, and a measure \(\mu\) that records which objects are
present.  The objects are observed through a strongly measurable map
\(\psi\) with values in a real Hilbert space \(Y\), and \(Y\) carries a
linear isometric involution \(\Jiso\) compatible with \(\psi\) in the
sense that \(\psi(Jx)=\Jiso\psi(x)\).  The involution \(\Jiso\) splits every
response into an even part and an odd part.  The odd part is given by
the orthogonal projection
\[
  \Pminus=\tfrac12\,(I_Y-\Jiso),
  \qquad
  \psim(x)=\Pminus\psi(x),
\]
and it changes sign under the duality: \(\psim(Jx)=-\psim(x)\).  In
particular \(\psim\) vanishes at every fixed point of \(J\).  When it
vanishes \emph{only} there, the readout \(\psim\) detects exactly
the displacement of an object from the fixed locus \(\Fix(J)\).

The obstruction is the positive operator
\[
  \AX=\int_X \psim(x)\,\psim(x)^*\,d\mu(x),
\]
where \(v^*\) denotes the functional \(h\mapsto\langle v,h\rangle\),
so that \(vv^*\) is the rank-one operator
\(h\mapsto\langle v,h\rangle v\).  When
\(\int_X\|\psim\|^2\,d\mu<\infty\), the operator \(\AX\) is positive
and trace-class, and
\[
  \trace\AX=\int_X\|\psim(x)\|^2\,d\mu(x).
\]
Suppose now that \(\AX\) is dominated by positive trace-class
operators whose traces tend to zero:
\[
  \AX\preceq B_n\quad(n\geq1),
  \qquad
  \trace B_n\to0,
\]
where \(\preceq\) is the Loewner order (\(A\preceq B\) means that
\(B-A\) is positive).  Monotonicity of the trace gives
\(0\leq\trace\AX\leq\trace B_n\) for every \(n\), so
\(\trace\AX=0\).  Hence \(\AX=0\) and \(\psim=0\) almost everywhere.
If \(\psim\) vanishes only on \(\Fix(J)\), then \(\mu\)-almost every
object is fixed by \(J\): the visible mass is \emph{confined} to the
fixed locus.  Conversely, if \(\psim=0\) almost everywhere then
\(\AX=0\) and \(B_n=0\) will do.  The equivalence is immediate; it is
recorded as \Cref{thm:duality_confinement:master-theorem}.

The argument is short, and none of its operator-theoretic ingredients
is new: the Loewner order \citep{Loewner1934,Bhatia1997}, trace
ideals \citep{Simon2005}, Markov's inequality \citep{Folland1999},
the arithmetic--geometric mean inequality
\citep{HardyLittlewoodPolya1952} and Douglas's factorization theorem
\citep{Douglas1966} are classical.  Nor does the theorem make
confinement easier to prove: by the equivalence just stated, producing
domination records is exactly as hard as proving
\(\int_X\|\psim\|^2\,d\mu=0\).  What the formulation offers is a
format for the case in which \(\AX\) is not accessible directly.  In
applications of the framework one computes, from other data, a
positive operator \(\Kminus\) on \(Y\)---for instance a carrier-side
quantity built, in finite dimensions, from an audit operator and a
probe \citep{TsiokosNeedleKiller2026}---and establishes, by a separate
argument, a bridge inequality \(\AX\preceq\Kminus+E\) with an explicit
positive defect \(E\).  Trace bounds on \(\Kminus\) and \(E\) then
control \(\trace\AX=\int_X\|\psim\|^2\,d\mu\), the odd energy of the
ledger; under quantitative separation this in turn bounds the mass at
a given distance from the fixed locus.
All of the substance lies in the bridge; the results of this paper
record what follows from it.

The paper therefore separates two things that are easy to conflate:
\begin{enumerate}
  \item the \emph{confinement theorem} and its variants, proved here,
    which convert domination records into confinement on the fixed
    locus; and
  \item a \emph{hypothesis} about which structures possess such
    records.  In the author's Six Birds framework
    \citep{TsiokosFoundationsI2026} this hypothesis is named
    Self-Dual Trace Confinement, written \(\GamSDTC\) (SDTC for
    short): in specified closures that carry a genuine involutive
    self-duality, trace-vanishing domination records for the odd
    readout are supplied by the structure.  SDTC is not proved here,
    and no result of this paper depends on it.  It is not a
    statement about all involutive ledgers: the first example in
    \Cref{sec:involutive_ledger}, taken with Lebesgue measure on a
    disc centred on the real axis, is an involutive ledger with a
    separating odd readout that is not confined, and hence has no
    such records.
\end{enumerate}
\Cref{sec:framework} explains the framework vocabulary in plain
terms; a reader interested only in the operator theory can skip it.

A guiding picture is the functional-equation symmetry of the Riemann
zeta function.  The map \(s\mapsto1-\bar s\) is an involution of the
complex plane whose fixed locus is the critical line
\(\operatorname{Re}s=\tfrac12\), and \(\operatorname{Re}s-\tfrac12\)
is the coordinate that changes sign.  A measure on the nontrivial
zeros, weighted by multiplicity, is confined to the fixed locus
exactly when the Riemann hypothesis holds.  The picture explains the
shape of the abstraction; no result about \(\zeta\) is used or
claimed in this paper.  A companion paper \citep{TsiokosRHviaSDTC2026}
studies this instance under explicitly stated hypotheses.

\paragraph{Results.}
Throughout, \((X,J,\mu,\psi,Y,\Jiso)\) is an involutive object ledger
(\Cref{def:duality_confinement:involutive-object-ledger}) with
\(\int_X\|\psim\|^2\,d\mu<\infty\).
\begin{enumerate}
  \item \emph{Trace identity}
    (\Cref{lem:duality_confinement:trace-identity}).  \(\AX\) is a
    positive trace-class operator with
    \(\trace\AX=\int_X\|\psim\|^2\,d\mu\).
  \item \emph{Direct confinement}
    (\Cref{thm:duality_confinement:separation-confinement,thm:duality_confinement:quantitative-confinement}).
    If \(\AX=0\) and \(\psim\) is separating, \(\mu\)-almost every
    point is fixed by \(J\).  If \(X\) is a metric space and
    \(\|\psim(x)\|\geq m(\varepsilon)>0\) whenever
    \(\dist(x,\Fix(J))\geq\varepsilon\), then the set of points at
    distance at least \(\varepsilon\) from \(\Fix(J)\) has outer measure
    at most \(\trace\AX/m(\varepsilon)^2\).
  \item \emph{Douglas factorization}
    (\Cref{thm:duality_confinement:douglas-domination}).  For bounded
    operators \(V,W\) with a common domain,
    \(V^*V\preceq W^*W\) if and only if \(V=TW\) for a contraction
    \(T\).  Applied to the analysis operator of \(\psim\), this
    describes exact domination of \(\AX\) as a factorization problem.
  \item \emph{Confinement theorem}
    (\Cref{thm:duality_confinement:master-theorem}).  If
    \(\AX\preceq B_n\) for positive trace-class \(B_n\) with
    \(\trace B_n\to0\), then \(\AX=0\), \(\psim=0\) almost
    everywhere, and, if \(\psim\) is separating, \(\mu\)-almost every
    point is fixed by \(J\).  Conversely, if \(\psim=0\) almost
    everywhere, the zero sequence is such a family of records.
  \item \emph{Finite windows and defect budgets}
    (\Cref{thm:duality_confinement:exhaustive-squeeze} and
    \Cref{prop:duality_confinement:optimized-trace-budget}).
    The same conclusion holds when the domination is supplied on
    finite windows: real Hilbert spaces \(Y_n\), bounded operators
    \(\iota_n:Y_n\to Y\), positive trace-class operators
    \(A_{X,n}\preceq B_n\) on \(Y_n\), and positive trace-class
    operators \(T_n\) on \(Y\), with
    \(\AX\preceq\iota_nA_{X,n}\iota_n^*+T_n\) and
    \(\trace(\iota_nB_n\iota_n^*)+\trace T_n\to0\).  When \(\AX\) is dominated by
    \((1+t)(K+E_{\mathrm{src}})+(1+t^{-1})E_{\mathrm{br}}\) for every
    \(t>0\), with \(K,E_{\mathrm{src}},E_{\mathrm{br}}\) positive
    trace-class, then
    \(\trace\AX\leq\bigl(\sqrt a+\sqrt b\bigr)^2\) with
    \(a=\trace(K+E_{\mathrm{src}})\), \(b=\trace E_{\mathrm{br}}\), and
    this value is the infimum over \(t>0\) of the scalar budgets.
\end{enumerate}
Every numbered result in this list, and the abstract squeeze of
\Cref{rem:duality_confinement:abstract-cone}, is formalized in Lean~4
with Mathlib for real Hilbert spaces of arbitrary dimension.  The formal
proofs take the trace of a positive operator to be the sum of its
diagonal entries in a Hilbert basis and the Gram operator to be a
Bochner integral in the operator norm.  Bochner integrability in the
trace class, the standard facts identifying the two traces, and the
application of Douglas factorization to the analysis operator are
checked only in the text.  \Cref{app:formalization} lists the
correspondence.

\paragraph{Organization.}
\Cref{sec:framework} explains the framework vocabulary.
\Cref{sec:involutive_ledger,sec:anti_invariant_ledger} define the
ledger, the odd readout, the operator \(\AX\) and separation.
\Cref{sec:direct_confinement} proves the direct confinement results,
\Cref{sec:domination} discusses domination and Douglas factorization,
\Cref{sec:master_theorem} proves the confinement theorem, and
\Cref{sec:exhaustive_squeeze_and_budgets} gives the finite-window and
defect-budget refinements.  \Cref{sec:scope,sec:discussion} discuss
scope and possible instances, and \Cref{sec:conclusion} concludes.
